(1) 线性互感磁链为 $\lambda_1=Li_1+Mi_2$、$\lambda_2=Mi_1+Li_2$。无源回路KVL为 $L\dot i_1+M\dot i_2+Ri_1=0$、$M\dot i_1+L\dot i_2+Ri_2=0$；互感的符号由绕组参考方向定义。

(2) 相加相减后 $(L+M)\dot i_++Ri_+=0$、$(L-M)\dot i_-+Ri_-=0$。因此 $i_\pm(t)=i_\pm(0)e^{-t/\tau_\pm}$，时间常数 $\boxed{\tau_\pm=(L\pm M)/R}$，均为正，因为 $|M|<L$。

(3) 初始两个正常坐标均为 $I_0/2$，故 $i_1(t)=I_0[e^{-t/\tau_+}+e^{-t/\tau_-}]/2$，$i_2(t)=I_0[e^{-t/\tau_+}-e^{-t/\tau_-}]/2$。若 $M>0$，有 $\tau_+>\tau_-$，$i_2$ 先按正方向感生；改变互感接线符号则反向。

(4) 磁能 $U_B=(Li_1^2+2Mi_1i_2+Li_2^2)/2=(L+M)i_+^2+(L-M)i_-^2>0$。对时间求导并利用KVL得 $\dot U_B=-R(i_1^2+i_2^2)\le0$。$M=0$ 时两个回路独立衰减；$M\to L^-$ 时反协同时间常数趋零，理想集中参数模型的超快过程不再可靠。
